% Kamp III.2, Lemma 6: recovering a fixed initial endpoint between two realizations.

\subsubsection{A fixed initial endpoint between two realizations}
\label{sec:p09-kamp-iii2-lemma6}

Kamp's Lemma~6 fixes one basic pattern and two realizations beginning at
the same point. Between their final points, a realization with that
fixed beginning exists exactly when some past realization of the pattern
ends there. The proof joins portions of two finite witness sequences
\cite[ch.~III, lemma~6, pp.~79--80]{Kamp1968}.

\begin{hbconvention}{def:p09-kamp-fixed-initial-standing}{One fixed pattern and its guarded past closure}
Fix distinct variable names \(x,y\), an integer \(n\geq1\), and
\(\phi\in\operatorname{Bp}_n(x,y)\). Its occurrence word is
\[
                 \sigma(\phi)=(A_0,B_1,A_1,\ldots,B_n,A_n)
\]
as in \cref{def:p09-kamp-basic-word}. All point and interval letters
have one fixed interpretation in the structure under consideration,
and repeated letters remain allowed. The integer \(n\) counts
intervals; each realizing sequence has exactly \(n+1\) points.

For a linear structure
\(\mathcal M=\langle T,<,(P_i^{\mathcal M})_i\rangle\), write
\(\phi(r,s)\) as a semantic abbreviation for
\(\mathcal M\models\phi[x\mapsto r,y\mapsto s]\), where \(r,s\in T\).
This does not introduce an additional relation symbol.
Define the one-variable formula
\[
 \phi^{\exists<}(y)\defeq
             \exists v\,(v<y\land\phi[v/x]),
\]
where \(v\) is fresh and the replacement concerns only free occurrences
of \(x\). Bound variables are renamed fresh beforehand if needed.
Write \(\phi^{\exists<}(s)\) for its truth with \(y=s\).
The superscript names this formula operation. It is the source's
guarded closure \((\exists x<y)\phi(x,y)\), with the bound initial
endpoint renamed to distinguish it from a separately fixed point.

By the basic-pattern realization equivalence already proved in
\cref{annex:p9.kamp.iii2.lemma1}, for \(r<s\) the assertion
\(\phi(r,s)\) means that the word has a realization from \(r\)
to \(s\). The assertion \(\phi^{\exists<}(s)\) means that some
realization ends at \(s\), with any strictly earlier initial point.
In particular this quantifier is not restricted by a fixed initial
point occurring in another premise
\cite[ch.~III, pp.~63--66, 79]{Kamp1968}.
\end{hbconvention}

\begin{hblemma}{thm:p09-kamp-iii2-lemma6}{proved=annex:p9.kamp.iii2.lemma6}{Kamp's fixed-initial-point lemma}
Let \(\mathcal M\) be any structure over a nonempty strict linear
order, and keep the fixed pattern above. Let
\[
                         a<b<c<d
\]
be four points of its domain. Suppose
\[
                         \phi(a,b),\qquad \phi(a,d).
\]
Then
\[
                     \phi(a,c)
                         \quad\Longleftrightarrow\quad
                                  \phi^{\exists<}(c).
\]
Both premises concern the same formula \(\phi\) with the same
predicate interpretation. The source's four underlined point symbols
\(\underline{t_1},\underline{t_2},\underline{t_3},\underline{t_4}\)
correspond respectively to \(a,b,c,d\); they are distinct from the
object-language variable names.

The statement explicitly assumes a linear order in the source, and
the proof applies throughout \(\mathsf{Lin}\). Completeness, density
and absence of endpoints are not hypotheses. The result therefore also
holds on \(\mathsf{Com}\), without transferring any earlier
complete-order equivalence to incomplete orders
\cite[ch.~III, lemma~6, pp.~79--80]{Kamp1968}.
\end{hblemma}

\begin{hbmethod}{def:p09-kamp-fixed-initial-splice}{Constructing the realization with the fixed beginning}
Suppose the right side of the lemma holds. The input to this finite
construction consists of a realizing sequence
\[
                  \mathbf v=(v_0,\ldots,v_n),\qquad v_n=c,
\]
and realizations \(\mathbf w^-\) from \(a\) to \(b\) and
\(\mathbf w^+\) from \(a\) to \(d\), supplied by the two premises.
Every sequence realizes the same fixed word \(\sigma(\phi)\).

If \(a\leq v_0\), take \(\mathbf w=\mathbf w^+\) and choose the
least index \(k\in\{0,\ldots,n-1\}\) satisfying
\[
                         w_k\leq v_k<w_{k+1}.
\]
If \(v_0<a\), take \(\mathbf w=\mathbf w^-\) and choose the
least index in that same set satisfying
\[
                         v_k\leq w_k<v_{k+1}.
\]
The annex proves existence of an eligible index in each case, and
proves the second splice also valid under its weak comparison
\(v_0\leq a\). Thus equality of the two initial points is covered
by either source argument.

Return the sequence
\[
                    (w_0,\ldots,w_k,v_{k+1},\ldots,v_n).
\]
It has exactly \(n+1\) points and realizes the same word from \(a\)
to \(c\). The index range includes both \(k=0\) and \(k=n-1\).
In the source's numbering from 1, its switch index is \(k+1\),
so this is exactly the printed range \(1,\ldots,n\).
The choice of the least eligible index merely makes the finite
construction definite; the proof works for every index meeting the
indicated source condition.

The method is a construction from supplied witnesses. It does not
claim to find arbitrary witnesses in a structure from a formula alone.
It uses only points in the two selected realizations, and it never
changes the number of pattern intervals
\cite[ch.~III, p.~80]{Kamp1968}.
\end{hbmethod}

\begin{hbexample}{def:p09-kamp-splice-overlap}{An old endpoint inside the new joining interval}
Take \(T=\{0,1,\ldots,8\}\) in its usual order and a two-interval
pattern with word \((A_0,B_1,A_1,B_2,A_2)\), using distinct
predicate letters interpreted by
\[
 \begin{gathered}
 A_0=\{0,1\},\qquad A_1=\{2,4\},\qquad A_2=\{3,6,8\},\\
 B_1=\{1,2,3\},\qquad B_2=T.
 \end{gathered}
\]
Here the displayed sets are the predicate interpretations.
Choose \((a,b,c,d)=(0,3,6,8)\).
The realizations \((0,2,3)\) and \((0,2,8)\) supply the two
premises, and \(\mathbf v=(1,4,6)\) supplies the guarded past
closure at 6.

The first case uses \(\mathbf w=(0,2,8)\). Its first eligible
switch is \(k=0\), since \(0\leq1<2\), and the output is
\((0,4,6)\). The new first gap \((0,4)\) contains both 1 and 2.
The old gap \((0,2)\) supplies \(B_1\) at 1, and the old gap
\((1,4)\) supplies it at 2 and 3. Thus the old point \(w_1=2\),
which is now inside the new gap, receives its interval predicate
from the other realization. Its old point predicate \(A_1\)
would not by itself supply that requirement. The annex checks the
whole example and the general joining interval.
\end{hbexample}

\begin{hbexample}{def:p09-kamp-splice-final-label}{The guarded past condition retains the final point label}
For a one-interval word \((A_0,B_1,A_1)\) with distinct predicate letters,
take \(T=\{0,1,2,3\}\) in its usual order. Interpret \(A_0,B_1\)
by \(T\), and \(A_1\) by \(\{1,3\}\).
The four points \(a=0,b=1,c=2,d=3\) satisfy both premises of
the lemma. Nevertheless \(\phi(0,2)\) and
\(\phi^{\exists<}(2)\) are both false, because their final
point must satisfy \(A_1\). The two premises therefore do not
assert that the pattern holds from \(a\) to every intermediate
point without the stated guarded past condition.
\end{hbexample}

\paragraph{The next source boundary.}
Lemma~7 begins on printed p.~80, immediately after this proof.
Its two parts transform guarded formulas with a universal interior
condition into finite alternatives of basic and elementary interval
conditions, with specified restrictions on the lengths of the resulting
disjuncts. Its statement continues on p.~81, and its further proof
is the next separate construction. The present result establishes
only the fixed-pattern assertion of Lemma~6
\cite[ch.~III, pp.~80--81]{Kamp1968}.
